\documentclass[../main.tex]{subfiles}
\begin{document}
\chapter{Fuerza en un conductor}
\img{images/04/campo.jpg}{Fuerza conductor que circula corriente
dentro de un campo magnético uniforme.}
\info{Fuerza sobre un conductor que circula corriente situado en un
campo magnético}{
  \begin{equation*}
    F= B \times l \times I
  \end{equation*}

  Siendo:

  $F$: Fuerza en Newtons $[N]$.

  $B$: Inducción Magnética en teslas $[T]$.

  $l$: Longitud del conductor en metros $[m]$.

  $I$: Corriente en Amperios. $[A]$.
}
\actividad{Resolver los siguientes problemas.}
\begin{enumerate}
  \item Un conductor es recorrido por una corriente de $I=10A$, tiene
    una longitud de $l=200m$ dentro de un campo magnético uniforme de
    inducción $B=9T$ y esta situado perpendicularmente a las líneas
    de fuerza. Calcular la fuerza que el campo magnético ejerce sobre
    el conductor.
  \item ¿Qué intensidad $I$ debe circular por un conductor que tiene
    $l=10cm$ de longitud dentro de un campo magnético de inducción
    $B=1400Gs$, si está cortando perpendicularmente a las líneas de
    fuerza del campo para que se ejerza sobre el conductor $F=0,5N$?
  \item Una bobina es recorrida por una corriente $I=15A$, tiene uno
    de sus lados de $l=10cm$ dentro de un campo magnético de
    $B=12000Gs$. Calcular la fuerza $F$ que el campo ejerce sobre el
    lado de la bobina de $N=100$ espiras.
  \item Un conductor es recorrido por $I=16A$ tiene una longitud de
    $l=20cm$ dentro de un campo magnético uniforme y situado
    perpendicularmente a la dirección del campo. Calcular la
    inducción siendo la fuerza de $F=5N$.
  \item Un conductor tiene una longitud de $l=25cm$ y está situado
    perpendicularmente a un campo de inducción $B=1,6T$, que ejerce
    una fuerza de $F=10N$. Calcular la intensidad de corriente $I$.
  \item Dibujar la dirección del las líneas de fuerza de campo
    magnético, luego determinar la fuerza sobre el conductor usando
    la regla de la mano izquierda o regla de Fleming.
    \begin{enumerate}
      \item.\subfile{../images/04/04-6-1}
      \item.\subfile{../images/04/04-6-2}
      \item.\subfile{../images/04/04-6-3.tex}
      \item.\subfile{../images/04/04-6-4.tex}
      \item.\subfile{../images/04/04-6-5.tex}
      \item.\subfile{../images/04/04-6-6.tex}
      \item.\subfile{../images/04/04-6-7.tex}
      \item.\subfile{../images/04/04-6-8.tex}
      \item.\subfile{../images/04/04-6-9.tex}
      \item.\subfile{../images/04/04-6-10.tex}
    \end{enumerate}
\end{enumerate}
\end{document}

